\documentclass[a4paper,11pt]{article}
\usepackage{exam-style-842-unit}
\examinehead{842 信息化概论 · 第 05 单元练习册}{练习二 参考答案与解析}
\begin{document}

\unitanswerscover{842 信息化概论 · 单元练习册}{第 05 单元参考答案}{练习二（云服务与虚拟化）的答案与解析；判断题均给改正表述，教材未展开的虚拟化分类、容器与 Hypervisor 均注明来源。}

\answerpart{一、单项选择题（10 题，每题 3 分，共 30 分）}

\q B。\begin{analysis}第二层是平台即服务（PaaS），向用户提供分布式操作系统、分布式数据库等基础软件，公司所有的开发都可以在这一层进行（教材 4.1.4，第 45 页）。A 位于最底层，提供服务器、存储、网络等资源；C 位于第三层，提供可直接使用的应用；D 教材未把 DaaS 列为三类之一。\end{analysis}

\q A。\begin{analysis}使用 IaaS 后用户可以购买服务器基础服务、使用即可；服务器可能是虚拟的、是庞大服务器集群的一部分，但用户可自由安装和定制任意操作系统（如 Windows 或 Linux），不必理会服务器硬件购买、托管、监控、维护等工作，只需专注于自己的网站或应用（教材 4.1.4，第 45 页）。B PaaS 提供基础软件；C SaaS 直接提供应用；D 私有云是部署方式而非服务模式。\end{analysis}

\q C。\begin{analysis}SaaS 是人们在生活中每天都要接触的一层，大多通过网页浏览器实现；用户消费的服务完全从桌面系统或智能手机中获取，如微信、微博、云存储、在线视频，国外则如 Netflix、Google Apps、Dropbox、iCloud（教材 4.1.4，第 46 页）。A、B 是资源与平台层；D 是部署方式。\end{analysis}

\q A。\begin{analysis}教材 4.1.6 把 EC2 列为 AWS 的计算类服务：EC2 是虚拟机实例，有标准型、大内存、高运算能力等多种类型虚拟主机，可自动按需伸缩（第 49 页）。虚拟机实例属于基础设施资源，即 IaaS；B PaaS 提供平台与基础软件；C SaaS 提供应用；D 是节能技术，与服务模式无关。\end{analysis}

\q A。\begin{analysis}教材 4.1.4 指出：亚马逊公司因为创立了 PaaS 商业模式而被《商业周刊》评为 2008 年的创新企业，亚马逊原本是出售图书的电子商务公司，后来发现自己的电子商务平台可以租用给别人（第 45 页）。B 微软、C Google、D 阿里云均与这一表述不符。\end{analysis}

\q B。\begin{analysis}Hypervisor（虚拟机监视器，VMM）是运行在物理硬件与客户操作系统之间的软件层，负责虚拟机的创建、调度与隔离（通行资料，教材 4.1.5 只讲“在软件中仿真计算机硬件”，未展开 Hypervisor，本册按通行资料补充）。A、C、D 描述的是应用、Web、存储之间的接口，与 Hypervisor 无关。\end{analysis}

\q B。\begin{analysis}容器与虚拟机最大的区别在于隔离层次：容器共享宿主机操作系统内核，以进程为隔离单位，不需要为每个容器安装完整的客户操作系统，因而启动快、镜像小、开销低；虚拟机则虚拟化硬件，每个虚拟机都包含完整的客户操作系统，隔离性更强但开销更大（通行资料，教材未展开）。A 与事实相反；C 容器正是以 Linux 内核的命名空间与 cgroups 为基础；D 容器属于操作系统级虚拟化。\end{analysis}

\q C。\begin{analysis}阿里云创立于 2009 年，是中国最大的云计算平台，服务范围覆盖全球 200 多个国家和地区（教材 4.1.6，第 50 页）。A 百度云、B 腾讯云、D 华为云虽然在市场上存在，但教材只把百度云、腾讯云与 Google 云、微软 Azure 并列为“其他云服务商”，并未称其为中国最大，也未提及华为云。\end{analysis}

\q A。\begin{analysis}教材 4.1.6 指出：AWS 于 2006 年推出，基础设施功能丰富，已为全球 190 个国家/地区内成百上千家企业提供支持，是世界上最大的公有云平台服务提供商（第 49—50 页）。B、C、D 虽然都是著名云平台，但教材没有给出“世界最大”的评价。\end{analysis}

\q B。\begin{analysis}教材 4.1.7 指出：云安全是通过网状的大量客户端监测网络中软件行为的异常，获取互联网中木马、恶意程序的最新信息，推送到服务器端进行自动分析和处理，再把病毒和木马的解决方案分发到每一个客户端（第 51 页）。A 因噎废食；C 只讲边界防护；D 只讲加密与离线保存，均不是云安全的完整含义。\end{analysis}

\answerpart{二、多项选择题（3 题，每题 4 分，共 12 分）}

\q ABC。\begin{analysis}A、B 正确：SaaS 大多通过网页浏览器实现，用户消费的服务完全从桌面系统或智能手机中获取（教材 4.1.4，第 46 页）。C 正确：微信、微博、云存储、在线视频都是教材列举的 SaaS 例子。D 错误：SaaS 用户直接使用应用，服务器与系统由服务商维护，用户无需自行维护操作系统和数据库。\end{analysis}

\q ABC。\begin{analysis}A 正确：虚拟化是在软件中仿真计算机硬件、以虚拟资源为用户提供服务的计算形式（教材 4.1.5，第 47 页）。B、C 正确：教材给出两种应用模式——把一台性能强大的服务器虚拟成多个独立的小服务器，以及把多台服务器虚拟成一台强大的服务器（第 47 页）。D 错误：教材明确“虚拟化是云计算的重要组成部分但不是全部”，二者不能画等号。\end{analysis}

\q ABC。\begin{analysis}教材 4.1.6 列出的 AWS 服务中，EC2 是虚拟机实例（计算类）、EBS 是可作为 EC2 实例持久性数据块级存储的存储服务、Elastic MapReduce 是基于 Hadoop 的 MapReduce 型分析服务（第 49 页）。D 微信属于 SaaS 应用，不是 AWS 提供的服务。\end{analysis}

\answerpart{三、判断题（6 题，每题 2 分，共 12 分）}

\q √。\begin{analysis}教材 4.1.4 指出：第二层就是平台即服务（PaaS），此服务层次向用户提供基础软件，如分布式操作系统、分布式数据库等基础服务（第 45 页）。\end{analysis}

\q ×。改正：使用 IaaS 的用户购买服务器基础服务、使用即可，不必自行购买、托管、监控和维护物理服务器，只需专注于自己的应用。\begin{analysis}教材 4.1.4 指出 IaaS 的好处在于：用户不用理会服务器硬件购买、托管、监控、维护等工作，只需专注于自己的网站（第 45 页）。“必须自己购买、托管和维护物理服务器”是使用 IaaS 之前的传统做法。\end{analysis}

\q ×。改正：容器共享宿主操作系统内核，不需要安装完整的客户操作系统；只有虚拟机才为每个实例安装完整的客户操作系统。\begin{analysis}该内容教材未展开，按通行资料：容器是操作系统级虚拟化，以进程为隔离单位；虚拟机是硬件级虚拟化，每台虚拟机都含完整客户操作系统。\end{analysis}

\q √。\begin{analysis}教材 4.1.6 指出：亚马逊网络服务（AWS）云于 2006 年推出；根据其页面介绍，AWS 已经为全球 190 个国家/地区内成百上千家企业提供了支持，是世界上最大的公有云平台服务提供商（第 49—50 页）。\end{analysis}

\q √。\begin{analysis}教材 4.1.6 指出：阿里云创立于 2009 年，是中国最大的云计算平台，服务范围覆盖全球 200 多个国家和地区（第 50 页）。\end{analysis}

\q ×。改正：云安全是通过大量客户端监测软件行为异常、由服务器端自动分析处理并把解决方案分发到客户端的技术；云上数据并非绝对安全，公有云还会增加安全风险。\begin{analysis}教材 4.1.7 给出云安全的定义（第 51 页）；教材 4.1.4 同时指出公有云“增加了安全风险”，且存在法律法规缺失、责任难以追究等缺陷（第 45 页）。因此“绝对安全”的说法错误。\end{analysis}

\answerpart{四、填空题（6 空，每空 1 分，共 6 分）}

\q SaaS。\begin{analysis}NIST 把云计算服务模式分为 IaaS、PaaS 和 SaaS 三大类或层次（教材 4.1.4，第 45 页）。\end{analysis}

\q 基础设施即服务。\begin{analysis}IaaS 即 Infrastructure-as-a-Service，位于云计算最底层，有时也叫 Hardware-as-a-Service（硬件即服务）（教材 4.1.4，第 45 页）。\end{analysis}

\q Platform as a Service（Platform-as-a-Service）。\begin{analysis}PaaS 即 Platform-as-a-Service，第二层，提供分布式操作系统、分布式数据库等基础软件（教材 4.1.4，第 45 页）。\end{analysis}

\q 虚拟化；服务器。\begin{analysis}教材 4.1.5 给出虚拟化定义：在软件中仿真计算机硬件，以虚拟资源为用户提供服务的计算形式（第 47 页）。把一台物理服务器划分为多个虚拟服务器属于服务器虚拟化（通行分类，教材只给出“一虚多”的应用模式）。\end{analysis}

\q 云存储。\begin{analysis}教材 4.1.7 指出：云存储是通过集群应用、网格技术或分布式文件系统等功能，将网络中大量各种不同类型的存储设备通过应用软件集合起来协同工作，共同对外提供数据存储和业务访问功能（第 51 页）。\end{analysis}

\answerpart{五、简答题（4 题，每题 6 分，共 24 分）}

\q 参考答案要点（教材 4.1.4，第 45—46 页）：
（1）IaaS（2 分）：位于最底层，向用户提供服务器、存储、网络等基本计算资源。用户可自由安装和定制任意操作系统与应用，负责操作系统及以上部分；硬件购买、托管、监控、维护由服务商负责。教材例子：EC2 虚拟机实例、EBS 存储。
（2）PaaS（2 分）：位于第二层，向用户提供分布式操作系统、分布式数据库等基础软件，以及开发和分发应用的解决方案；用户专注于在此层开发自己的应用，不必管理底层硬件和基础软件。教材例子：亚马逊平台服务（教材称亚马逊创立了 PaaS 商业模式）。
（3）SaaS（2 分）：位于第三层，向用户直接提供应用软件，用户通过网页浏览器或智能终端使用，无需安装、维护软件与服务器。教材例子：微信、微博、云存储、在线视频、Google Apps、Dropbox、iCloud。
（4）可补充：三层自下而上抽象程度提高，用户直接控制的范围递减、开箱即用的程度递增。

\q 参考答案要点（教材 4.1.5，第 47 页）：
（1）定义（2 分）：虚拟化是一种在软件中仿真计算机硬件、以虚拟资源为用户提供服务的计算形式，旨在合理调配计算机资源，使其更高效地提供服务；它打破了应用系统各硬件间的物理划分，实现架构的动态化，实现物理资源的集中管理和使用。
（2）两种应用模式（2 分）：一是把一台性能强大的服务器虚拟成多个独立的小服务器，服务不同的用户；二是把多台服务器虚拟成一台强大的服务器，完成特定的功能。两种模式的核心都是统一管理、动态分配资源、提高资源利用率。
（3）最大好处（2 分）：增强系统的弹性和灵活性，降低成本、改进服务、提高资源利用效率。可补充：虚拟化是云计算最重要的核心技术之一，为云计算服务提供基础架构层面的支撑，但虚拟化不等于云计算。

\q 参考答案要点（教材未展开，按通行资料作答）：
（1）隔离层次不同（2 分）：虚拟机在硬件层虚拟化，每个虚拟机拥有独立的客户操作系统和虚拟硬件，隔离性强；容器在操作系统层虚拟化，共享宿主机内核，以进程（命名空间、cgroups）为隔离单位，隔离性相对弱。
（2）启动速度与资源开销不同（2 分）：虚拟机启动通常需数十秒甚至更久，镜像大、占用内存与存储多；容器可秒级启动，镜像小、开销低，同一宿主机能运行更多实例。
（3）适用场景不同（2 分）：虚拟机适合运行不同操作系统、需要强隔离的多租户与遗留应用场景；容器适合微服务、持续集成与持续交付、需要快速弹性伸缩的应用部署。实际生产中二者常结合使用（虚拟机中运行容器）。
（4）说明：教材 4.1.5 只讲虚拟化的定义与两种应用模式，未展开 Hypervisor、容器等分类，本小题按通行资料作答即可。

\q 参考答案要点（教材 4.1.6，第 49—51 页）：
（1）AWS（2 分）：亚马逊网络服务，2006 年推出，是世界上最大的公有云平台服务提供商，已为全球 190 个国家/地区内成百上千家企业提供支持。主要服务包括 EC2 虚拟机实例、EBS 数据块存储、SimpleDB、RDS、SQS、CloudFront、VPC、CloudWatch、Elastic MapReduce 等，并提供开发工具、SDK、文档、社区与技术支持。
（2）阿里云（2 分）：创立于 2009 年，是中国最大的云计算平台，服务范围覆盖全球 200 多个国家和地区；提供政务、游戏、金融、电商、移动、医疗、多媒体、物联网等行业解决方案；曾抵御峰值 453.8Gbit/s 的 DDoS 攻击，其分布式计算平台 ODPS 以 377 秒完成 100TB 数据排序；在全球多地部署绿色数据中心。
（3）其他主流平台（2 分）：教材还提到 Google 云、微软 Azure 云、百度云、腾讯云等也颇为著名。答题时应注意教材未提及华为云，若补充应说明其为通行资料而非教材表述。

\answerpart{六、综合题（1 题，共 16 分）}

\q 参考答案要点：
（1）服务模式判断（6 分，每项 1.5 分，判断与理由各半）：
第一，农产品电商网站——IaaS。企业自建网站，需要自主安装和定制操作系统、Web 服务器与数据库，只希望省去硬件购买、托管、监控与维护，符合教材 IaaS 的典型场景（4.1.4，第 45 页）。
第二，ERP 与办公系统——SaaS。属于标准化成熟软件，员工通过浏览器或客户端直接使用，无需自行维护服务器与软件（4.1.4，第 46 页）；也可答 PaaS，但须说明理由。
第三，育种数据分析平台——PaaS。平台向团队提供分布式操作系统、分布式数据库等基础软件与开发环境，各课题组在这一层开发分析应用；若强调直接租用弹性算力与存储，也可答 IaaS，理由对应“提供服务器、存储、网络等基础资源”即可。
在线课程——SaaS。教材 4.1.7 把 MOOC（大规模开放在线课程）列为云教育的典型应用，属于直接使用的应用服务。
（2）虚拟化的作用与服务器虚拟化的好处（4 分）：
作用（2 分）：虚拟化把物理资源抽象为资源池，实现统一管理、动态分配，提高资源利用率，增强系统的弹性和灵活性，降低成本和改进服务（4.1.5，第 47 页）；它是云计算提供基础架构层面支撑的最重要核心技术之一。
服务器虚拟化的好处（2 分，答出 3 点即可）：把一台物理服务器划分为多个隔离的虚拟机，提高 CPU、内存、存储的利用率；减少物理服务器数量与机房能耗，降低成本；虚拟机可快速创建、克隆、迁移和备份，便于弹性伸缩与故障恢复（通行资料，可结合教材“一虚多”模式作答）。
（3）部署模式组合与安全隐私（6 分）：
部署组合（2 分）：面向公众的电商网站与在线课程宜用公有云，可发挥规模经济效益、按需弹性扩展（4.1.4，第 45 页）；内部 ERP 与育种数据涉及企业核心数据，宜用私有云，部署在防火墙内、只为企业内部服务，安全性与服务质量更可控（第 46 页）；两者组合即混合云（第 46 页）。
安全与隐私风险（2 分）：公有云增加了安全风险，存在法律法规不健全、责任难以追究、计费标准不统一、运营商数量不多、更换服务商时数据迁移困难等缺陷（4.1.4，第 45 页）；云上还存在多租户隔离不足、账号被盗、数据泄露与滥用等风险。
应对措施（2 分，答出 4 点即可，属通行整理）：选择合规、信誉好的云服务商并签订明确的数据归属与责任条款；对敏感数据加密存储与传输，做好访问控制与最小权限管理；采用多副本与异地备份，定期演练恢复；开启审计日志与安全监测，利用云安全机制及时发现木马与异常行为（4.1.7，第 51 页）；对公开数据进行脱敏处理，并事先约定退出云服务商时的数据迁移方案。

\end{document}
